\documentclass[main.tex]{subfiles}
\begin{document}
\chapter{Bogomolov's Theorem on Curves of Bounded Genus}

%% All commands are in macros.sty

\section{The Boundedness Theorem of Bogomolov}

References:

\begin{itemize}
\item [1] Bogomolov's original article ``Families of curves on a surface of general type'' (1977) which is scanned on \chref{https://mathoverflow.net/questions/52419/rational-curves-on-varieties-of-general-type}{mathoverflow}.
\item [2] The S\'{e}minaire N. Bourbaki article ``Courbes de genre g\'{e}om\'{e}trique born\'{e} sur une surface de type g\'{e}n\'{e}ral'' (1979) by Mireille Martin-Deschamps. The article is hosted \chref{https://www.numdam.org/book-part/SB_1977-1978__20__233_0/}{here} by Numdam. Warning, there are some errors in this text.
\item [3] Miles Reid's article ``Bogomolov's theorem $c_1^2 \le 4 c_2$'' (1977) found \chref{https://mreid.warwick.ac.uk/Reprint/Bogomolov.pdf}{here} which proves that $c_1(X)^2 \le 4 c_2(X)$ for a smooth projective general type surface. Note many of the statements in this talk do not appear in this article but the proofs are all essentially there.
\end{itemize}

\begin{theorem}[Bogomolov]
Let $X$ be a smooth projective surface of general type with $c_1^2 > c_2$. Then curves of fixed genus on $X$ form a bounded family.
\end{theorem}

\begin{rmk}
Note if $f_t : C_t \to X$ is a family of curves whose image has nontrivial deformation then the family covers $X$. Hence its normal bundle has nonnegative degree. Therefore,
\[ 2g(C_t) - 2 = (K_X + C_t) \cdot C_t \ge K_X \cdot C_t \]
which means if we fix the degree then the degree of the image curves on the canonical model of $X$ is bounded. This implies that the covering families of fixed genus are bounded. Hence the content of Bogomolov's theorem is about rigid curves.
\end{rmk}

\begin{cor}
Since $X$ cannot have a covering family of genus $0$ or $1$ curves there can be at most finitely many such curves on $X$.
\end{cor}

\begin{defn}
We say that a vector bundle $E$ is \textit{big} if $\struct{}(1)$ on $\P(E)$ is big. Equivalently
\[ h^0(X, S^m E) \in O(m^{\dim{X} + \rank{E} - 1}) \]
\end{defn}

Really the theorem has two parts:
\begin{enumerate}
\item if $\Omega_X$ is big then curves of fixed genus are bounded
\item if $K_X$ is big and $c_1^2 > c_2$ then $\Omega_X$ is big.
\end{enumerate}

\begin{rmk}
In fact, to conclude $X$ has only finitely many genus $0$ or $1$ curves we only require $H^0(X, S^m \Omega_X) \neq 0$ for some $m > 0$.
\end{rmk}

The first step follows the same strategy as the other hyperbolicity results: we know curves on $X$ lift to $\P(\Omega_X)$ -- the first step in the jet tower -- and they come in two types: regular curves -- those not landing in a base locus $Z \subset \P(\Omega_X)$ -- and irregular curves -- those landing in $Z$. The former are bounded by looking at their degree against $\struct{\P(E)}(1)$. The latter are bounded because they are leaves of a certain foliation.
\par
We explain bigness of $\Omega_X$ first. A Riemann-Roch calculation gives
\[ \chi(X, S^m \Omega_X) = \tfrac{1}{3!} (c_1^2 - c_2) m^3 + O(m^2) \]
so if $c_1^2 > c_2$ this grows maximally and we win if $h^2(X, S^m \Omega_X)$ has growth bounded by $O(m^2)$. In fact, it is zero for $m \ge 3$ by Bogomolov's vanishing theorem.

\begin{theorem}[Bogomolov]
Let $X$ be a smooth projective surface of general type. Then $H^0(X, S^m \Omega_X \ot \omega_X^{-q}) = 0$ whenever $2q > m$.
\end{theorem}

\begin{rmk}
Hence $H^2(X, S^m \Omega_X \ot \omega_X^r) = H^0(X, S^m \Omega_X \ot \omega_X^{1 - m - r})^{\vee}$ vanishes for $2(r + m - 1) > m$. In particular for $r = 0$ and $m \ge 3$.
\end{rmk}

Bogomolov's vanishing theorem follows from analysis of the stability of the cotangent bundle.

\subsection{Notions of Stability}

Say that a class $D \in \NS(X)$ is \textit{positive} if $D^2 > 0$ and $D \cdot H > 0$ for all ample $H$. Say that a line bundle $A$ is $\Q$-\textit{effective} if $A^{\ot m}$ has a section for some $m > 0$.

\begin{defn}
Let $E$ be a rank $2$ vector bundle on a surface $X$. Suppose there is a sequence
\[ 0 \to A \to E \to B \ot I_Z \to 0 \]
with $A,B$ line bundles and $Z$ a finite subscheme. Setting $\Delta = A \ot B^{-1}$
\begin{enumerate}
\item if $c_1(\Delta) \in N(X)^+$ is a positive class we say $E$ is \textit{Bogomolov-unstable}
\item if $\Delta$ is $\Q$-effective and either $\Delta$ is nontorsion or $Z \neq \empty$ then we say $E$ is \textit{weakly Bogomolov unstable}
\end{enumerate}
\end{defn}

\begin{defn}
We say that a rank $2$ vector bundle $E$ on a surface $X$ is \textit{Mumford unstable} if there exists a nonzero section of $(S^{2m} E) \ot (\det{E})^{-m}$ which attains a zero somewhere.
\end{defn}

\begin{theorem}[Bogomolov-Reid]
The following implications hold for a rank $2$ vector bundle on a smooth projective surface $X$,
\begin{center}
\begin{tikzcd}
c_1(E)^2 > 4 c_2(E) \arrow[r, equals] & \text{Bogomolov unstable} \arrow[d]  \arrow[r] & \text{weakly Bogomolov unstable} \arrow[d,equals]
\\
& \mu_H\text{-unstable for all } H \arrow[d] & \text{M-unstable}
\\
& \mu_H\text{-unstable for some } H
\end{tikzcd}
\end{center}
\end{theorem}

\begin{example}
$\mu_H$-unstable for all $H$ does not imply weakly Bogomolov unstable. Consider $X = \Bl_p \P^2$ and $E$ be the exceptional divisor. Then let
\[ \E = \struct{X}(E) \oplus \struct{X} \]
Clearly, $\E$ is $\mu_H$-unstable for all $H$ since $H \cdot E > 0$. However, I claim that $\E$ is not weakly Bogomolov unstable. Otherwise we would have $A + B = E$ and $P = A - B$ positive or zero. Thus $A = \tfrac{1}{2}(P + E)$ so $H \cdot A > 0$ for all $H$. Also at least one of $E - A$ or $-A$ is effective by the sequence. Since $A \cdot H > 0$ the second cannot hold but $\struct{}(A) \to \E$ is a saturated subsheaf so $A = E$ so $B = 0$ but $A - B$ is neither positive nor zero giving a contradiction.
\bigskip\\
% REVIEW: Should "$\mu_H$-stable" here be "$\mu_H$-unstable"? The diagram above shows the implication goes from Bogomolov unstable to $\mu_H$-unstable, so the claim should be that $\mu_H$-unstable does not imply M-unstable.
In fact, $\mu_H$-unstable for all $H$ does not imply Mumford unstable. Indeed, we need to find $X$ with a divisor $D$ such that $D \cdot H > 0$ for all $H$ but $D$ is not $\Q$-effective. In this case
\[ \E = \struct{X}(D) \oplus \struct{} \]
gives an example. Indeed, we can just take $D$ on $\P^1 \times E$ where $E$ is an elliptic curve to be $\pi_1^* ([0]) + \pi_2^* ([p] - [o])$ where $p \in E$ is a nontorsion-point. If we want $D$ to not even be numerically equivalent to a $\Q$-effective divisor we can do this on a more complicated ruled surface, see Lazarsfeld positivity I example 1.5.1 for a pseudoeffective but not $\R$-effective class.
\end{example}


The relevance of these properties is the following lemma.

\begin{lemma}
Let $E$ be an M-stable rank $2$ vector bundle with $\Q$-effective and nontorsion determinant. Then $H^0(X, S^m E \ot (\det{E})^{-q}) = 0$ for $2 q > m$.
\end{lemma}

\begin{proof}
Let $L = \det{E}$ and suppose there is a section
\[ s \in H^0(X, S^m E \ot L^{-q}) \]
choose a positive integer $k$ large enough that there is a nonzero section
\[ w \in H^0(X, L^{k(2q - m)}) \]
which is possible because $2q > m$ and $L$ is $\Q$-effective. Then using the symmetric algebra structure
\[ s^{2k} \cdot w \in H^0(X, S^{2km} E \ot L^{-km}) \]
is a nonzero section. Since $L$ is nontorsion $w$ has a zero somewhere. This shows that $E$ is M-unstable.
\end{proof}

Therefore, to prove Bogomolov's vanishing result, we just need to show that $\Omega_X$ is M-stable. To do this, we will show $\Omega_X$ is not weakly Bogomolov unstable.


\subsection{Stability of the Cotangent Bundle}

\begin{theorem}
Let $X$ be a smooth projective general type surface. Then $\Omega_X$ is not weakly Bogomolov unstable.
\end{theorem}

\begin{example}
It is possible for $\Omega_X$ to be $\mu_H$-unstable for all ample $H$ even if $X$ is general type.  Consider $X = \Bl (C_1 \times C_2)$. Then
\[ K_X = K_{C_1} + K_{C_2} - E \]
and if $C_1, C_2$ are two general curves every ample is of the form
\[ H = H_1 + H_2 - c E \]
where $H_i$ are the pullbacks of ample classes on $C_i$. Then,
\[ \mu_H(\Omega_X) = \tfrac{1}{2} ((2g_2 - 2) \deg{H_1} + (2g_1 - 2) \deg{H_2} - c) \]
but we have subsheaves $\pi^* \Omega_{C_i}$ with slope
\[ \mu_H(\pi^* \Omega_{C_i}) = (2g_i - 2) \deg{H_i'} \]
where $i'$ is the opposite of $i$.
The larger of these two numbers destabilizes $\Omega_X$. However, we do see that $\Omega_X$ is not Bogomolov unstable.
\end{example}

The proof of the stability theorem follows immediately from the following two results.

\begin{lemma}
Let $E$ weakly Bogomolov unstable. If $\det{E}$ is big then $E$ contains a big line bundle.
\end{lemma}

\begin{proof}
Consider the sequence
\[ 0 \to A \to E \to B \ot I_Z \to 0 \]
where by assumption $A \ot B^{-1}$ is $\Q$-effective. Therefore, for some $m > 0$ we have $B^m \embed A^m$. However, $A \ot B \cong \det{E}$ is big. The injection
\[ A^m \ot B^m \embed A^{2m} \]
then proves that $A$ is also big.
\end{proof}

\begin{theorem}[Bogomolov]
If $X$ is a smooth projective general type surface and $\L \embed \Omega_X^1$ is a line bundle subsheaf then $h^0(X, \L^m) \in O(m)$.
\end{theorem}

This will follow immediately from a lemma of Castelnuovo and de Franchis:

\begin{lemma}
Let $X$ be a smooth projective surface. Suppose $\L \subset \Omega_X^1$ is a line bundle with $h^0(X, \L) \ge 2$ then there exists a fiber space $f : X \to C$ over a curve so that
\[ f^* \Omega_C \subset \L \subset f^* \Omega_C(D) \]
where $D \subset C$ is an effective divisor supported on the locus where $f$ has a nonreduced fiber. % REVIEW: The $g(C)$ notation for genus conflicts with the map $g : X' \to C$ used in the proof. Consider writing $\mathrm{genus}(C) \ge 2$ or renaming one of them.
Moreover $\Omega_C \subset f_* \L$ is an isomorphism so $g(C) \ge 2$.
\end{lemma}

\begin{proof}[Proof of Theorem]
We suppose that $\L \embed \Omega_X^1$ is an invertible subsheaf. If $h^0(X, \L^{\ot n}) \le 1$ for all $n$ then we are done. Otherwise, there is some $n > 0$ such that $h^0(X, \L^{\ot n}) \ge 2$. In this case, by passing to a cyclic cover we may assume that $h^0(X, \L) \ge 2$. Then the lemma implies $h^0(X, \L^{\ot m}) \in O(m)$.
\end{proof}

\begin{proof}[Proof of Lemma]
Therefore, there are two independent $1$-forms $\omega_1, \omega_2 \in H^0(X, \L) \subset H^0(X, \Omega^1_X)$ such that $\omega_1 \wedge \omega_2 = 0$ because they lie in the same $1$-dimensional subspace at the generic point $\L_{\eta} \subset \Omega_{X, \eta}^1$. Therefore, there is a rational function $f : X \rat \P^1$ so that $\omega_1 = f \omega_2$. Resolving the map $\pi : X' \to X$ we have $f' : X' \to \P^1$. The Stein factorization produces $g : X' \to C \to \P^1$. Since $\omega_i$ are global $1$-forms
\[ 0 = \d{\omega_1} = \d{f} \wedge \omega_2 \]
and therefore $\d{f}$ as a rational $1$-form lives inside $\L_{\eta}$ because it is parallel to $\omega_2$. If $t$ is the coordinate on $\P^1$ then $\d{f} = f^* \d{t}$ so we see the inclusions
\[ f^* \Omega_{\P^1} \subset g^* \Omega_{C} \subset \pi^* \L \subset \Omega_{X'} \]
by checking at the generic point as submodules of $\Omega_{X',\eta}$. This works because at the generic point $f^* \Omega_{\P^1}$ and $g^* \Omega_C$ are equal because $C \to \P^1$ is generically \etale. Now it suffices to check that
\begin{enumerate}
\item $\Omega_C = g_* \L$
\item $\L \subset g^* \Omega_C(D)$
\end{enumerate}
For the first, the inclusion $\Omega_C \subset g_* \L$ allows us to view a section of $g_* \L$ as a meromorphic $1$-form on $C$. However, because of the maps
\[ g^* \Omega_C \to g^* g_* \L \to \L \subset \Omega_{X'} \]
this meromorphic $1$-form becomes regular after pulling back to $X'$. Now we use the following somewhat surprising lemma.
\end{proof}

\begin{lemma}
Let $f : X \to Y$ be a dominant map of smooth proper varieties. Suppose $\omega \in \Omega_{Y,\eta}^k$ is a meromorphic $k$-form such that $f^* \omega$ extends to a global regular $k$-form $f^* \omega \in H^0(X, \Omega_X^k)$ on $X$. Then $\omega \in H^0(Y, \Omega_Y^k)$ is regular.
\end{lemma}

\begin{rmk}
This is false in positive characteristic even for separable maps! The issue is failure of tameness. It also fails for symmetric forms $S^m \Omega_X$ for $m > 1$. Indeed, $z \mapsto z^2$ can regularize $z^{-1} (\d{z})^2$.
\end{rmk}

% REVIEW: The argument below completes the proof of the Castelnuovo--de Franchis lemma (items (1) and (2)) but sits outside the \begin{proof}...\end{proof} environment, which ended earlier. Consider restructuring so this material is inside the proof, or adding a separate proof environment.
To complete the proof we need to show that $\L \subset g^* \Omega_C(D)$ for $D$ the divisor of points over which $f$ has a multiple fiber (counted each with multiplicity $1$). At the generic point of a component of a multiple fiber, the map $X \to C$ is formally equivalent to
\[ \CC\dbrac{t} \to \CC \dbrac{x,y} \quad t \mapsto x^n \]
where $n$ is the multiplicity of the component. Then
\[ \d{t} \mapsto n x^{n-1} \d{x} \]
meaning that
\[ \L \subset (g^* \Omega_C)(D(g)) \subset \Omega_X \]
where $D(g)$ is the \textit{ramification divisor}
% REVIEW: The summation variable $C$ here (components of fibers) conflicts with the curve $C$ that is the base of the fibration $g : X' \to C$. Consider renaming one of them.
\[ D(g) = \sum_{C \subset X_t} (n_C - 1) [C] \]
where $n_C$ is the multiplicity of a component of the fiber $X_t$. In fact, $(g^* \Omega_C)(D(g))$ is exactly the saturation of $g^* \Omega_C$ inside $\Omega_X$. But by definition,
\[ g^* \struct{C}(D) = \struct{X}(\sum_{C \subset X_t} n_C [C]) \supset \struct{X}(D(g)) \]
so
\[ \L \subset g^* (\Omega_C(D)) \]
and we win.

\subsection{Boundedness}

Proof of boundedness: given a nonconstant map $f : C \to X$ there is a lift $\wt{f} : C \to \P(\Omega_X)$ defined at the generic point by
\[ \d{f} : f^* \Omega_X \to \Omega_C \]
and then extending using properness. The image of $\d{f}$ is a line subbundle identified with $\Omega_C(-D)$ where $D$ is the ramification of $f$. By definition the universal quotient
\[ \pi^* \Omega_X \onto \struct{\P}(1) \]
pulls back to
\[ \wt{f}^* \pi^* \Omega_X \to \wt{f}^* \struct{\P}(1) \]
equals
\[ f^* \Omega_X \to \Omega_C(-D) \subset \Omega_C \]

Since $\struct{\P}(1)$ is big, the locus $Z \subset \P$ where $\struct{\P}(m)$ does not separate points or tangent vectors for all $m$ is a proper subset. What I mean is that $\struct{\P}(m)$ for $m \gg 0$ defines a rational map
\[ t_m : \P(\Omega_X) \rat \P^N  \]
which is an embedding away from $Z$. Now we split the curves on $X$ into two types
\begin{enumerate}
\item regular curves: those such that $\wt{f} : C \to \P$ meets $Z$ properly
\item irregular curves: those such that $\wt{f} : C \to \P$ lands inside $Z$.
\end{enumerate}

For regular curves, the composition
\[ \wt{f} : C \to \P(\Omega_X) \to \P^N \]
is nonconstant. This structure gives maps
\[ \wt{f}^* t_m^* \struct{\P^N}(1) \embed \wt{f}^* \struct{\P}(m) \embed (\Omega_C)^{\ot m} \]
Hence
\[ \deg{(t_m \circ \wt{f})} \le m (2g - 2) \]
Since $m$ can be chosen once and for all, we see that the regular curves of fixed genus map to $\P^N$ with bounded degree and hence are bounded.
\par
For the irregular curves, they must lift into some irreducible component $Z_i \subset Z \subset \P(\Omega_X)$. Let $g : Y \to Z_i$ be a resolution. Either $Y \to X$ has only a single curve in its image in which case we are done or $Y \to X$ is surjective. Pulling back the sequence
% REVIEW: The kernel is called $K$ in the exact sequence but $Q$ afterwards -- unify notation (perhaps $Q$ throughout, to avoid conflict with canonical divisor $K$).
\[ 0 \to Q \to \pi^* \Omega_X \to \struct{\P}(1) \to 0 \]
to $Y$ gives
\[ g^* Q \to g^* \pi^* \Omega_X \to g^* \Omega_{\P} \to \Omega_Y \]
Since $Y \to X$ is dominant and generically finite $g^* Q \to \Omega_Y$. The kernel of its dual defines a foliation $\F \subset \T_Y$ which is integrable because it has rank $1$. Furthermore, any curve $f : C \to X$ lifting in the canonical way to $Y$ must satisfy
\begin{center}
\begin{tikzcd}
0 \arrow[d]
\\
\wt{f}^* Q \arrow[d]
\\
f^* \Omega_X \arrow[r] \arrow[d] & \wt{f}^* \Omega_{\P} \arrow[r] & \wt{f}^* \Omega_Y \arrow[d]
\\
\wt{f}^* \struct{\P}(1) \arrow[d] \arrow[rr] & & \Omega_C
\\
0
\end{tikzcd}
\end{center}
so $\wt{f}^* Q \to \Omega_C$ is zero meaning $C$ is tangent to the foliation $\F$ on $Y$. We finish by applying the following theorem of Jouanolou

\begin{theorem}
Let $\F \subset \Omega_X^1$ be an algebraic foliation on a smooth projective variety over $\CC$. Either there are only finitely many algebraic hypersurfaces tangent to $\F$ or $\F$ is induced by the fibers of an algebraic contraction $X \to Y$.
\end{theorem}

Either way, we get boundedness for the curves on $Y$ tangent to the foliation. Since every irregular curve lifts to a curve tangent to a foliation on one of the finitely many choices for $Y$, this completes the proof.

\ifSubfilesClassLoaded{%
  \bibliographystyle{alpha}
  \bibliography{refs}
}{}

\end{document}
